\chapter{Stanislav Smirnov}\label{ch2010stanislavsmirnov:stanislav-smirnov-fields-medal-2010-a-comprehensive-report-on-his-work-and-impact}

\begin{center}\rule{0.5\linewidth}{0.5pt}\end{center}

\textbf{Prepared:} July 2026 \textbf{Subject:} Stanislav Smirnov (b. September 3, 1970, Leningrad, USSR) \textbf{Position:} Professor, University of Geneva; Head, Laboratory of Mathematical Physics, Chebyshev Laboratory, St.~Petersburg State University \textbf{Doctoral Advisor:} Nikolai Makarov (Caltech) \textbf{Fields Medal Citation:} \emph{``For the proof of conformal invariance of percolation and the planar Ising model in statistical physics.''}

\begin{center}\rule{0.5\linewidth}{0.5pt}\end{center}

\begin{prereqs} Probability theory (random walks, Markov chains, conformal invariance), complex analysis (holomorphic functions, discrete complex analysis), statistical mechanics (Ising model, percolation). \end{prereqs}

\begin{objectives} Prove the conformal invariance of critical percolation on the triangular lattice; describe the discrete complex analysis approach to lattice models \cite{smirnov2010}; explain Cardy's formula and its proof. \end{objectives}

\section{Problem History}\label{ch2010stanislavsmirnov:problem-history}

\textbf{Critical percolation} on a lattice is a fundamental model of randomness in statistical mechanics. Each edge (or site) of the lattice is independently ``open'' with probability $p$ and ``closed'' with probability $1-p$. The critical probability $p_c$ is the threshold at which an infinite connected cluster of open edges first appears. For site percolation on the triangular lattice, $p_c = 1/2$.

At criticality ($p = p_c$), the model is expected to be \emph{conformally invariant}: the law of the interfaces between open and closed clusters, in the scaling limit, is invariant under conformal transformations of the domain. This means that crossing probabilities---the probability that there exists an open path connecting two boundary arcs---depend only on the conformal type of the domain, not on its specific geometry.

The conformal invariance of critical percolation had been conjectured by physicists (using conformal field theory) since the 1980s. John Cardy (1984) derived a formula for the crossing probability in a rectangular domain:

\[
\Pi(a,b) = \frac{\Gamma(2/3)}{\Gamma(4/3)\Gamma(1/3)} \, {}_2F_1(1/3, 2/3; 4/3; r),
\]

where $r$ is the aspect ratio of the rectangle and ${}_2F_1$ is the hypergeometric function. This \textbf{Cardy's formula} was a prediction; proving it required showing that the scaling limit exists and is conformally invariant.

The \textbf{planar Ising model} is another central model. At the critical temperature, the model exhibits a second-order phase transition, and the interfaces between spin-up and spin-down clusters are expected to converge to a conformally invariant random curve. Smirnov proved conformal invariance for the Ising model \emph{before} the percolation case, using similar discrete complex analysis methods.

The triangular lattice is special for percolation because the observables are exactly discrete holomorphic on this lattice, but not on others. This exact holomorphicity is what makes the proof work for the triangular lattice; extending to other lattices (e.g., square, hexagonal) requires different techniques.

\section{How Was It Solved}\label{ch2010stanislavsmirnov:how-was-it-solved}

\textbf{Discrete Complex Analysis.}

Smirnov's key innovation is the method of \textbf{discrete complex analysis}: studying discrete holomorphic functions on the lattice and showing they converge to holomorphic functions in the scaling limit. A function $f$ on the vertices of a planar graph is \emph{discrete holomorphic} if its value at each vertex is the average of its neighbors---or more generally, if it satisfies a discrete analogue of the Cauchy--Riemann equations.

For critical percolation on the triangular lattice, Smirnov defined the \textbf{cardinality observable}: for a simply connected domain $D$ with boundary arcs $a, b, c, d$ (in cyclic order), define

\[
F(z) = \sum_{\gamma} e^{i \theta(\gamma)} \mathbb{P}(\gamma \text{ connects } a \text{ to } c),
\]

where the sum is over all open paths $\gamma$ from $a$ to $c$ passing through $z$, and $\theta(\gamma)$ is the winding angle of $\gamma$. Smirnov proved that $F$ is \emph{discrete holomorphic} on the triangular lattice.

\textbf{Conformal Invariance of Crossing Probabilities.}

By showing that $F$ converges to a holomorphic function in the scaling limit (using the Riemann mapping theorem and normal families arguments), Smirnov proved that the crossing probability

\[
\mathbb{P}(\text{open path from } a \text{ to } c)
\]

is a conformal invariant of the domain. This implies Cardy's formula: in a rectangle of aspect ratio $r$, the crossing probability depends only on $r$, and must equal the hypergeometric function derived by Cardy.

The proof has three main steps:
\begin{enumerate}
  \item \textbf{Discrete holomorphicity}: Show that the cardinality observable $F$ satisfies the discrete Cauchy--Riemann equations on the triangular lattice.
  \item \textbf{Convergence}: Prove that $F$ converges to a holomorphic function $f$ in the scaling limit, using the fact that the lattice mesh $\to 0$.
  \item \textbf{Boundary values}: Determine the boundary values of $f$ and use the Riemann mapping theorem to identify $f$ uniquely, yielding Cardy's formula.
\end{enumerate}

\textbf{The Ising Model.}

For the critical Ising model on a planar lattice, Smirnov defined a similar observable using the spin-spin correlation function. He showed that this observable is discrete holomorphic, and that its scaling limit is holomorphic. This established conformal invariance of the Ising model interfaces.

Smirnov's work on the Ising model (with Chelkak) also proved the existence of \emph{conformal invariance for the fermionic observables}, which are the scaling limits of certain discrete holomorphic functions related to the Ising model. The fermionic observable is a solution to the \emph{discrete holomorphicity condition}

\[
\sum_{w \sim z} (f(w) - f(z)) \, \tau(z,w) = 0,
\]

where $\tau(z,w)$ is a complex weight depending on the lattice geometry. For the Ising model on the square lattice, the weight is $\tau(z,w) = e^{i\pi \omega(z,w)/3}$ where $\omega(z,w)$ is related to the edge weight in the FK representation.

\textbf{Cardy's Formula:}

Cardy's formula states that the crossing probability in a rectangle of aspect ratio $r$ is:

\[
\Pi(r) = \frac{\Gamma(2/3)}{\Gamma(4/3)\Gamma(1/3)} \, {}_2F_1(1/3, 2/3; 4/3; r),
\]

where ${}_2F_1(a,b;c;z)$ is the Gaussian hypergeometric function. For $r=1$ (a square), $\Pi(1) = 1/2$, which is the probability that a percolation exploration path crosses from one side of the square to the opposite side. The formula was verified numerically before Smirnov's proof.

\section{Impact of This Discovery}\label{ch2010stanislavsmirnov:impact-of-this-discovery}

Smirnov's proof of conformal invariance for percolation and the Ising model was the first rigorous demonstration that discrete statistical mechanics models can have conformally invariant scaling limits. This confirmed 30 years of physical predictions and established \emph{discrete complex analysis} as a powerful method in probability theory.

The identification of the percolation scaling limit as SLE$_6$ and the Ising scaling limit as SLE$_3$ (proved by Smirnov and collaborators) connected discrete lattice models to the SLE framework developed by Schramm, Lawler, and Werner. This connection is now a central paradigm in the field.

In mathematical physics, conformal invariance provides a tool for computing critical exponents exactly. The exact values of the critical exponents for percolation and the Ising model follow from the conformal invariance, without the need for detailed lattice calculations. These exact values serve as benchmarks for numerical simulations and experimental measurements in condensed matter physics.

The discrete complex analysis approach has been applied to other models, including the \emph{loop-erased random walk} (SLE$_2$), the \emph{uniform spanning tree} (SLE$_8$), and the \emph{dimer model}. It has also inspired developments in \emph{discrete differential geometry} and the theory of \emph{discrete conformal maps}.

\section{Modern Views and Improvements}\label{ch2010stanislavsmirnov:modern-views-and-improvements}

The discrete complex analysis program has been extended to \emph{multiply connected domains} and to \emph{random planar maps}. The conformal invariance of the \emph{Fortuin--Kasteleyn--Renyaard} model (random cluster model) has been established in several cases using Smirnov's methods.

The \emph{Conformal Loop Ensemble} (CLE$_6$) has been identified as the scaling limit of the percolation exploration path, and the CLE framework provides a natural setting for studying multiple interfaces in percolation and related models.

Open problems include: proving the conformal invariance of percolation on \emph{non-triangular} lattices (where the observable is not exactly holomorphic), understanding the scaling limit of the \emph{critical Ising model on random planar maps}, and extending the methods to \emph{three-dimensional} lattice models (where conformal invariance does not hold).

The \emph{random planar map} is a graph embedded in the sphere, considered up to orientation-preserving homeomorphisms. The scaling limit of uniformly random planar maps is the Brownian map, proved by Le Gall and Miermont. The critical Ising model on random planar maps has been studied by several authors, and the conformal invariance of the scaling limit is expected to follow from the conformal invariance of the Ising model on planar lattices combined with the conformal invariance of the Brownian map. The key difficulty is that the random planar map is not a smooth surface, and the discrete complex analysis methods must be adapted to the random geometry.

The \emph{Ising model on random planar maps} has been studied using the bijection between planar maps and planar maps with a root face (the Tutte bijection) combined with the FK representation. The partition function of the Ising model on a random planar map can be expressed as a sum over configurations of the FK model, and the conformal invariance of the scaling limit follows from the conformal invariance of the FK model on the planar map. The exact critical exponents for the Ising model on random planar maps have been computed using conformal field theory, and they differ from the exponents on a fixed lattice due to the randomness of the geometry.

The \emph{universality conjecture}---that the scaling limit of a wide class of critical lattice models is conformally invariant and is described by SLE---remains a central open problem. Smirnov's methods provide the main tool for attacking this conjecture in two dimensions.

The \emph{Fortuin--Kasteleyn} (FK) random cluster model is a generalization of percolation and the Ising model. Each edge is either ``open'' with probability $p$ or ``closed'' with probability $1-p$, and the weight of a configuration is $p^{\text{open}} (1-p)^{\text{closed}} q^{\text{clusters}}$, where $q$ is the cluster weight parameter. For $q=1$, the FK model reduces to ordinary percolation; for $q=2$, it is the FK representation of the Ising model. The conformal invariance of the FK model for $q=2$ was proved by Smirnov and Chelkak using a discrete holomorphic observable, and the result has been extended to other values of $q$ by Duminil-Copin and others. The FK model is central to the study of random cluster boundaries and their scaling limits.

Recent work by Duminil-Copin and Smirnov established the conformal invariance of the self-avoiding walk on the hexagonal lattice (SLE$_{8/3}$), confirming a long-standing conjecture and extending the reach of Smirnov's discrete complex analysis methods. Duminil-Copin went on to win the Fields Medal in 2022 for related work on the Ising model.

Recent developments in the \emph{quantum gravity} program, connecting random planar maps to Liouville quantum gravity, have renewed interest in the scaling limits of lattice models on random geometries. Smirnov's discrete complex analysis techniques provide a starting point for this program.

The \emph{discrete Gaussian free field} (dGFF) is a random function on the vertices of a lattice that is the discrete analogue of the continuum Gaussian free field. The dGFF has covariance given by the discrete Green's function, and its scaling limit is the continuum GFF. The connection between the dGFF and SLE is through the \emph{discrete holomorphic observable}: the exponential of the dGFF, properly normalized, gives a discrete holomorphic function whose scaling limit is a continuum holomorphic function. This connection, developed by Smirnov and others, provides a unified framework for studying lattice models and their conformally invariant scaling limits.

The \emph{exploration path} is a random curve constructed from a percolation configuration by following the boundary between open and closed clusters. For critical percolation on the triangular lattice, the exploration path is a random curve that starts at a boundary point and grows into the interior of the domain, always keeping an open cluster on its left and a closed cluster on its right. Smirnov proved that the exploration path converges to SLE$_6$ in the scaling limit, and this convergence is the key step in the proof of conformal invariance. The exploration path is constructed using a \emph{discrete harmonic function} (the cardinality observable), and the convergence to SLE$_6$ follows from the convergence of the discrete harmonic function to a continuum harmonic function.

The \emph{Fortuin--Kasteleyn--Renyaard} (FK) representation of the Ising model represents the spin configuration as a random cluster model with parameter $q = 2$. The FK representation has the advantage that the spin-spin correlation function can be expressed as a two-point function of the random cluster model, and the conformal invariance of the FK model implies the conformal invariance of the Ising model. Smirnov and Smirnov--Chelkak proved conformal invariance of the FK Ising model by constructing a discrete holomorphic observable for the FK model, analogous to the cardinality observable for percolation. The FK Ising model is the $q = 2$ case of the Fortuin--Kasteleyn random cluster model, and its scaling limit is SLE$_{16/3}$ (for the cluster boundaries) and SLE$_3$ (for the spin interfaces).

The \emph{universality} of critical exponents---the fact that they depend only on the symmetry class of the model, not on microscopic details---is a central prediction of renormalization group theory. Smirnov's proof of conformal invariance provides a rigorous foundation for universality in two dimensions, and the exact values of critical exponents serve as benchmarks for numerical simulations and experimental measurements in condensed matter physics.

\begin{center}\rule{0.5\linewidth}{0.5pt}\end{center}

\section*{Summary}
\addcontentsline{toc}{section}{Summary}

Stanislav Smirnov proved the conformal invariance of critical percolation on the triangular lattice and the planar Ising model by developing the method of discrete complex analysis. The key idea is to define discrete holomorphic observables on the lattice and show they converge to holomorphic functions in the scaling limit, yielding conformal invariance of crossing probabilities and correlation functions. His proof of Cardy's formula and the identification of the scaling limits as SLE$_6$ and SLE$_3$ established a rigorous connection between discrete statistical mechanics and conformal field theory.

\begin{center}\rule{0.5\linewidth}{0.5pt}\end{center}
